\section*{Proyecto de Vida --- Heidy Quiroga Colque}

\subsection*{Quién soy}
Mi nombre es Heidy Quiroga Colque, tengo 17 años y vivo en San Julián, Santa Cruz, Bolivia. Estudio en el Bachillerato Técnico Humanístico (BTH) del Módulo Tecnológico Productivo San Julián, en la especialidad de Sistemas Informáticos, y provengo de la Unidad Educativa Progreso (6to A), donde curso mi último año de secundaria. Me considero una estudiante con futuro y con ambición. Me gusta la tecnología y hacer soluciones con la computadora que ayuden a las personas. Por eso quiero seguir la carrera de Sistemas.

\subsection*{Mis metas a corto, mediano y largo plazo}
A corto plazo quiero graduarme del colegio con buenas notas y obtener mi título de Bachiller Técnico Humanístico en Sistemas Informáticos, aprovechando cada materia técnica para llegar bien preparada a la siguiente etapa. A mediano plazo quiero entrar a la universidad y estudiar Ingeniería de Sistemas con dedicación, aprender más programación y diseño de páginas, y empezar a trabajar en proyectos de tecnología que me den experiencia y me permitan ayudar a mi familia. A largo plazo quiero trabajar como ingeniera en Sistemas, ayudar a la sociedad con soluciones tecnológicas que beneficien a mi comunidad, y tener una vida estable en todo sentido, sobre todo económicamente.

\subsection*{Lo que este proyecto significa para mí}
Este proyecto lo hicimos para apoyar a nuestro Módulo. Antes la inscripción era con filas largas y papeles que se perdían. Ahora todo es más ordenado y más justo para cerca de 400 postulantes cada año. Este proyecto me ayuda a crecer: aprendí a hacer pantallas fáciles para el celular, a avisar cuando un dato está mal y a cuidar los avances. Preguntar a 60 postulantes y hablar con 5 funcionarios me ayudó a explicar con palabras simples. Trabajar con Sharon me enseñó a trabajar en equipo. Desde recuperar la contraseña en menos de 2 minutos hasta repartir los paralelos por colegio, todo es experiencia real para mi futura carrera de Sistemas.
